%% ch08 --- asyncio dla tych, którzy znają Task
%%
%% Zalążek, wygenerowany przez tools/gen_stubs.py z tools/chapters.json.
%% Poniższy opis jest kontraktem tego rozdziału: co ma uzasadnić, co ma dać
%% czytelnikowi i czego NIE ma powtarzać za tomami towarzyszącymi. Napisanie
%% rozdziału oznacza usunięcie zalążkowego bloku (tego, który zaczyna się po
%% linii \chapter) i zastąpienie go rozdziałem -- po czym ten plik nie jest
%% już generowany.

\chapter{asyncio dla tych, którzy znają Task}
\label{ch:asyncio}

\chapterstub{%
\textit{[Opis do przetłumaczenia --- patrz wydanie angielskie.]}

Argument: the reader already has the mental model and most of it transfers;
the part that does not is where the incidents come from. The map: a
coroutine against a \code{Task} (cold against hot \dash{} a coroutine does
nothing until it is scheduled, a C\# \code{Task} is already running); an
event loop with no \code{SynchronizationContext} (no \code{ConfigureAwait},
no deadlock by \code{.Result}, but one blocking call freezes everything);
\code{TaskGroup} against \code{Task.WhenAll}, and \code{gather} with its
orphaning semantics; \code{CancellationToken} against \code{Task.cancel()}
and \code{CancelledError} (cancellation is an exception delivered at an
\code{await}, not a flag you poll); \code{asyncio.timeout} against
\code{CancelAfter}; threads against processes under the GIL, and where
\code{to\_thread} and a process pool belong; \code{httpx.AsyncClient}
against \code{HttpClient} (lifetime, pooling, and a default timeout that
exists); the sync boundary, \code{asyncio.run}, and the absence of an
\code{async Main}. Payoff: a translation table from the TPL to asyncio, and
the three incidents \dash{} a blocking call, an orphaned task, a swallowed
\code{CancelledError}. Measurement: experiment E4, the one-blocking-call
degradation reproduced in the TPL-against-asyncio framing, free. Five
exercises, one of them porting a C\# \code{WhenAll} plus
\code{CancellationToken} snippet. Overlap rule, and it is the strictest in
the book: the LangChain book's chapters 3 and 4 own the event loop,
cancellation and the measurements in depth. This chapter is the translation
view. It points there in one sentence and never repeats a section.

\medskip
\textbf{Budżet słów:} 3000.\\
\textbf{Wydanie:} v0.2.\\
\textbf{Pomiar:} E4.
}
